To better understand the contribution of each component in the proposed
objective, we conduct two controlled ablation experiments under the same
training and evaluation protocol.

\paragraph{Experimental Setup.}
All ablation experiments are performed using the same \texttt{nnUNet}
configuration, data splits, preprocessing, and training schedule as in
the main experiments. All models are evaluated using voxel-level and lesion-level metrics,
including Dice, HD95, lesion-wise F1, small lesion recall, and false
positive characteristics, to provide a comprehensive assessment of model
behavior in the sparse lesion regime.

\subsection*{Ablation I: Contribution of CAT and MIL Terms}

To isolate the effect of each component, we compare three variants:
(i) nnUNet-CAT, which incorporates only the structure-level
component-adaptive term,
(ii) nnUNet-MIL, which incorporates only the lesion-level supervision
term, and
(iii) nnUNet-CATMIL, which combines both terms into a unified objective.

Table~\ref{tab:ablation_cat_mil} reports the mean performance across
5-fold cross-validation. The results reveal a clear trade-off between the
two terms. nnUNet-CAT achieves the highest lesion-wise F1 (0.8386) but
lower small lesion recall (0.8230), indicating stronger structural
consistency but limited sensitivity to small lesions. In contrast,
nnUNet-MIL improves small lesion recall (0.8712) but degrades boundary
accuracy (HD95 = 9.2375 mm) and lesion-wise F1 (0.7492), reflecting less
precise segmentation. The combined nnUNet-CATMIL achieves the most balanced performance. It
obtains the best Dice (0.7834) and lowest HD95 (7.9817 mm), while also
achieving the highest small lesion recall (0.8730) and the lowest FP blob
fraction (0.2094). This indicates that combining CAT and MIL mitigates
their individual limitations, improving both detection of small lesions
and overall segmentation quality without increasing false positives.

\begin{table}[H]
\centering
\caption{Ablation study on the contribution of CAT and MIL terms.}
\label{tab:ablation_cat_mil}
\begin{tabular}{lcccccc}
\hline
Model & Dice & HD95 $\downarrow$ & Lesion F1 & Small Lesion Recall & FP Vol & FP Blob \\
\hline
nnUNet-CAT    & 0.7812 & 8.6844 & 0.8386 & 0.8230 & 1528.72 & 0.2277 \\
nnUNet-MIL    & 0.7805 & 9.2375 & 0.7492 & 0.8712 & 1539.77 & 0.2195 \\
nnUNet-CATMIL & \textbf{0.7834} & \textbf{7.9817} & 0.7571 & \textbf{0.8730} & 1536.88 & \textbf{0.2094} \\
\hline
\end{tabular}
\end{table}

\subsection*{Ablation II: Sensitivity to Loss Weights}

We analyze the impact of different weighting coefficients
$\lambda_{\text{CAT}}$ and $\lambda_{\text{MIL}}$ on model performance.
The results in Table~\ref{tab:ablation_lambda} reveal a clear trade-off
between segmentation accuracy and lesion-level detection.

\begin{table}[H]
\centering
\caption{Sensitivity analysis of $\lambda_{\text{CAT}}$ and $\lambda_{\text{MIL}}$.}
\label{tab:ablation_lambda}
\begin{tabular}{lccccc}
\hline
Model & Dice & HD95 $\downarrow$ & Lesion F1 & Small Recall & FP Vol $\downarrow$ \\
\hline
CATMIL0101 & 0.7675 & 9.3881 & 0.7720 & 0.8514 & 2292.17 \\
CATMIL0102 & \textbf{0.7954} & \textbf{5.9771} & 0.7452 & 0.8711 & 1345.83 \\
CATMIL0201 & \underline{0.7921} & \underline{6.7964} & \textbf{0.7754} & \textbf{0.8784} & \textbf{1233.92} \\
CATMIL0202 & 0.7663 & 9.4507 & 0.7632 & \underline{0.8699} & \underline{1729.42} \\
\hline
\end{tabular}
\end{table}

Lower weights (CATMIL0102) lead to the best global segmentation
performance, achieving the highest Dice (0.7954) and lowest HD95
(5.9771 mm), indicating improved overlap and boundary accuracy. However,
this comes with reduced lesion-wise F1 (0.7452), suggesting weaker
instance-level consistency. In contrast, higher weights (CATMIL0201) improve lesion-level behavior,
achieving the highest lesion F1 (0.7754), highest small lesion recall
(0.8784), and lowest FP volume (1233.92 mm$^3$). This indicates stronger
detection of small and sparse lesions with better control of false
positives, albeit with slightly reduced Dice and boundary accuracy. These results confirm that $\lambda_{\text{CAT}}$ and
$\lambda_{\text{MIL}}$ control the balance between voxel-level accuracy
and lesion-level sensitivity. Moderate weighting provides a balanced
solution, while higher weights emphasize detection and lower weights
favor precise segmentation.